\documentclass{article}
\usepackage[utf8]{inputenc}
\usepackage{amsmath,amssymb}
\usepackage[most]{tcolorbox}
\usepackage{geometry}
\geometry{margin=2.5cm}

\newcommand{\step}[1]{\par\noindent\hangindent=3.5em\hangafter=1%
  \makebox[3.5em][l]{\textbf{#1.}}}
\newcommand{\steptag}[1]{\hfill{\small\textsf{[#1]}}}

\newtcolorbox{proofbox}[1]{
  colback=gray!5, colframe=gray!60,
  fonttitle=\bfseries, title={#1},
  breakable, enhanced,
  left=4pt, right=4pt, top=4pt, bottom=4pt,
  before skip=10pt, after skip=10pt
}

\begin{document}

\begin{proofbox}{After first fixing one global witness package W\_RF suppli...}
\step{1} After first fixing one global witness package W\_RF supplied by lem-routef-raw-factor-setting-formation, for every input (H,Phi,eta) to which that formation result applies, fix one def-routef-raw-factor-setting datum S over that same W\_RF supplied by the same result; for every integer n >= 1 and every X in M\_n(S.B), writing the fields of (W\_RF,S) as the unqualified symbols below: Raw factor-map norms: with C\_V, C\_T, rho\_T from (1.1), for 0 <= eta <= rho\_T, every amplification satisfies (1-C\_V*eta)*||X|| <= ||tilde-Delta\_n X|| <= (1+C\_V*eta)*||X|| and max\{||tilde-Delta||\_cb, ||tilde-Upsilon||\_cb\} <= 1+C\_T*eta. \steptag{validated} \\[6pt]
\step{1.1} Factor branch: in the ambient finite-dimensional raw-factor setup, let A=Im(tilde-Phi), let v:B->A be the extended isomorphism furnished below, and define tilde-Delta=v (with codomain included in B(H)); then for every n and X, (1-C\_V*eta)||X|| <= ||tilde-Delta\_n X|| <= (1+C\_V*eta)||X||, hence ||tilde-Delta||\_cb <= 1+C\_V*eta. \steptag{validated} \\[6pt]
\step{1.1.1} Radius and domain: (1.1) defines bar-C\_E=max\{1,C\_E\}, C\_V=bar-C\_E*C\_A and rho\_T as a minimum containing rho\_AI and epsilon\_E/C\_A. Thus 0<=eta<=rho\_T implies eta<=rho\_AI and, by lem-routef-ai-defect-linearization, epsilon\_AI(eta)<=C\_A*eta<=epsilon\_E. In the ambient finite-dimensional setup A=Im(tilde-Phi) is finite-dimensional. \steptag{validated} \\[6pt]
\step{1.1.1.1} \textsc{assume}: Universal-context unpacking, not an added premise: root node 1 fixes W\_RF supplied by lem-routef-raw-factor-setting-formation and then fixes an arbitrary input (H,Phi,eta) to which that result applies and a datum S supplied by it. Hence the quantified domain and conclusion of lem-routef-raw-factor-setting-formation give that H is nonzero and finite-dimensional, Phi:B(H)->B(H) is UCP, 0<=eta<=rho\_id\textasciicircum{}corr, ||Phi\textasciicircum{}2-Phi||\_cb<=eta, S.tilde-Phi is the displayed functional-calculus map, and S.A=Im(S.tilde-Phi). In particular, no literal identification H=C\textasciicircum{}d is assumed or needed. \steptag{validated} \\[6pt]
\step{1.1.1.1.1} The amended binding prefix of root node 1 universally fixes precisely a W\_RF, an input (H,Phi,eta) in the domain of lem-routef-raw-factor-setting-formation, and a supplied setting datum S. Universally instantiate that validated external at this input: its hypotheses/domain give H nonzero finite-dimensional, Phi UCP, 0<=eta<=rho\_id\textasciicircum{}corr, and ||Phi\textasciicircum{}2-Phi||\_cb<=eta, while its conclusion says that the fields of S include the canonical functional-calculus tilde-Phi and A=Im(tilde-Phi). Thus every fact asserted in node 1.1.1.1 is inherited from the amended root and the allowed external; none is introduced as a new hypothesis. \steptag{validated} \\[6pt]
\step{1.1.1.1.1.1} Instantiate lem-routef-raw-factor-setting-formation at the global W\_RF and the arbitrary input (H,Phi,eta) fixed in root node 1, and select the datum S that root node 1 fixes from its existential conclusion. The external's domain supplies: H is nonzero finite-dimensional, Phi is UCP, 0<=eta<=rho\_id\textasciicircum{}corr, and ||Phi\textasciicircum{}2-Phi||\_cb<=eta. Its conclusion identifies the selected datum's tilde-Phi with the canonical functional-calculus formula and its A with Im(tilde-Phi). \steptag{validated} \\[4pt]
\step{1.1.1.2} Equation (1.1) uses rho\_AI:=eta\_A and defines rho\_T as a minimum containing rho\_AI and epsilon\_E/C\_A. Hence 0<=eta<=rho\_T gives 0<=eta<=eta\_A and C\_A*eta<=epsilon\_E. Applying lem-routef-ai-defect-linearization to the nonzero H, UCP Phi, and cb-defect bound fixed in 1.1.1.1 yields epsilon\_AI(eta)<=C\_A*eta<=epsilon\_E (and supplies the extended epsilon\_AI(eta)-C*-algebra structure used by sibling 1.1.2). \steptag{validated} \\[4pt]
\step{1.1.1.3} Because H=C\textasciicircum{}d with d<infinity, B(H) has finite dimension d\textasciicircum{}2. The range A=Im(tilde-Phi) of the linear map tilde-Phi:B(H)->B(H) is a linear subspace of B(H), so dim(A)<=d\textasciicircum{}2 and A is finite-dimensional. \steptag{validated} \\[4pt]
\step{1.1.1.4} Correction superseding the strict inequality in node 1.1.1.3: let d=dim\_C(H), which is a finite positive integer by 1.1.1.1. Then dim\_C(B(H))=d\textasciicircum{}2. Since tilde-Phi:B(H)->B(H) is linear and A=Im(tilde-Phi), the rank bound for a linear map gives dim\_C(A)=rank(tilde-Phi)<=dim\_C(B(H))=d\textasciicircum{}2<infinity. Equality can occur (for example tilde-Phi=I), so no strict inequality is asserted or needed. Therefore A is finite-dimensional. \steptag{validated} \\[4pt]
\step{1.1.2} By lem-routef-ai-defect-linearization, A=Im(tilde-Phi), with its inherited operator-space structure, is an extended epsilon\_AI(eta)-C*-algebra. By lem-thmainext-conditional there are a finite-dimensional C*-algebra B and an extended delta-isomorphism v:B->A with delta=C\_E*epsilon\_AI(eta); moreover delta<=C\_E*C\_A*eta<=bar-C\_E*C\_A*eta=C\_V*eta. \steptag{validated} \\[4pt]
\step{1.1.3} By def-extended-delta-inclusion, every amplification v\_n of the extended delta-isomorphism satisfies (1-delta)||X||<=||v\_nX||<=(1+delta)||X||. Since 0<=delta<=C\_V*eta, these bounds imply (1-C\_V*eta)||X||<=||v\_nX||<=(1+C\_V*eta)||X||. Under the raw-factor definition tilde-Delta=v (followed by the isometric inclusion A subset B(H)); hence these are the claimed amplification bounds, and their upper half gives ||tilde-Delta||\_cb<=1+C\_V*eta. \steptag{validated} \\[4pt]
\step{1.2} Inverse-composition branch: with tilde-Upsilon:=v\textasciicircum{}(-1) tilde-Phi, one has ||tilde-Upsilon||\_cb <= (1+C\_theta*eta)/(1-C\_V*eta). \steptag{validated} \\[6pt]
\step{1.2.1} A UCP map is completely contractive: for every n put Psi=id\_\{M\_n\} tensor Phi. Complete positivity makes id\_\{M\_2\} tensor Psi positive; applying it to the positive block matrix [[X\textasciicircum{}*X,X\textasciicircum{}*],[X,I]] and taking the Schur complement gives Psi(X)\textasciicircum{}*Psi(X)<=Psi(X\textasciicircum{}*X). Positivity and unitality give 0<=Psi(X\textasciicircum{}*X)<=||X||\textasciicircum{}2 I, hence ||Psi(X)||<=||X||. Therefore ||Phi||\_cb<=1. \steptag{validated} \\[4pt]
\step{1.2.2} Since rho\_T<=rho\_theta, lem-routef-functional-calculus-closeness applies and gives ||tilde-Phi-Phi||\_cb<=C\_theta*eta. The preceding complete contractivity and the triangle inequality yield ||tilde-Phi||\_cb<=1+C\_theta*eta. \steptag{validated} \\[6pt]
\step{1.2.2.1} In the notation feeding (1.1), rho\_theta=1/8, while rho\_T is a minimum containing rho\_theta; hence eta<=rho\_T implies 0<=eta<=1/8. Therefore lem-routef-functional-calculus-closeness gives ||tilde-Phi-Phi||\_cb<=C\_theta*eta. Node 1.2.1 gives ||Phi||\_cb<=1, so the cb-norm triangle inequality yields ||tilde-Phi||\_cb<=||Phi||\_cb+||tilde-Phi-Phi||\_cb<=1+C\_theta*eta. \steptag{validated} \\[4pt]
\step{1.2.3} Reapplying the extended-isomorphism estimate from lem-routef-ai-defect-linearization, lem-thmainext-conditional, and def-extended-delta-inclusion gives ||v\_nX||>=(1-C\_V*eta)||X|| for every n. Because (1.1) gives C\_V*eta<=1/4, v\_n is bijective and ||v\_n\textasciicircum{}(-1)||<=1/(1-C\_V*eta). Thus tilde-Upsilon\_n=v\_n\textasciicircum{}(-1) tilde-Phi\_n has norm at most (1+C\_theta*eta)/(1-C\_V*eta), uniformly in n; taking the supremum gives the stated cb bound. \steptag{validated} \\[6pt]
\step{1.2.3.1} The map v supplied by lem-thmainext-conditional is an extended isomorphism, so v is bijective and every amplification v\_n=id\_\{M\_n\} tensor v is bijective with inverse id\_\{M\_n\} tensor v\textasciicircum{}(-1). The extended-inclusion lower bound, enlarged as in the factor branch, is ||v\_nX||>=(1-C\_V*eta)||X||. Since (1.1) gives C\_V*eta<=1/4, it follows that ||v\_n\textasciicircum{}(-1)||<=1/(1-C\_V*eta). \steptag{validated} \\[4pt]
\step{1.2.3.2} By definition tilde-Upsilon=v\textasciicircum{}(-1) tilde-Phi, hence tilde-Upsilon\_n=v\_n\textasciicircum{}(-1) tilde-Phi\_n. Combining the uniform inverse bound with ||tilde-Phi\_n||<=1+C\_theta*eta and submultiplicativity yields ||tilde-Upsilon\_n||<=(1+C\_theta*eta)/(1-C\_V*eta) for every n; taking the supremum over n gives the claimed cb bound. \steptag{validated} \\[6pt]
\step{1.2.3.2.1} Validated node 1.2.3.1 gives, for every n, that v\_n is bijective and ||v\_n\textasciicircum{}(-1)|| <= (1-C\_V*eta)\textasciicircum{}(-1), with 1-C\_V*eta>0. Validated node 1.2.2.1 gives ||tilde-Phi||\_cb <= 1+C\_theta*eta, hence ||tilde-Phi\_n|| <= 1+C\_theta*eta for every n. Since tilde-Upsilon=v\textasciicircum{}(-1) composed with tilde-Phi and tilde-Phi\_n has range in M\_n(A), amplification commutes with composition, so tilde-Upsilon\_n=v\_n\textasciicircum{}(-1) composed with tilde-Phi\_n. Therefore submultiplicativity gives ||tilde-Upsilon\_n|| <= ||v\_n\textasciicircum{}(-1)|| ||tilde-Phi\_n|| <= (1+C\_theta*eta)/(1-C\_V*eta) for every n. Taking the supremum over n yields ||tilde-Upsilon||\_cb <= (1+C\_theta*eta)/(1-C\_V*eta). \steptag{validated} \\[4pt]
\step{1.3} Scalar assembly: for 0<=eta<=rho\_T, the guards in (1.1) imply C\_theta*eta,C\_V*eta<=1/4 and therefore (1+C\_theta*eta)/(1-C\_V*eta) <= 1+(C\_theta+3*C\_V)*eta=1+C\_T*eta; also C\_V<=C\_T, so the two branch bounds imply max\{||tilde-Delta||\_cb,||tilde-Upsilon||\_cb\}<=1+C\_T*eta, while the factor branch already gives the asserted amplification inequalities. \steptag{validated} \\[6pt]
\step{1.3.1} Put a=C\_theta*eta and b=C\_V*eta. The two guards [4(1+C\_theta)]\textasciicircum{}(-1) and [4(1+C\_V)]\textasciicircum{}(-1) in rho\_T give 0<=a<=C\_theta/[4(1+C\_theta)]<=1/4 and 0<=b<=C\_V/[4(1+C\_V)]<=1/4, so 1-b>0. \steptag{validated} \\[4pt]
\step{1.3.2} Direct algebra gives [1+a+3b]-(1+a)/(1-b)=b*(2-a-3b)/(1-b). Because 0<=a,b<=1/4, one has 2-a-3b>=1, hence the difference is nonnegative and (1+C\_theta*eta)/(1-C\_V*eta)<=(1+C\_theta*eta+3*C\_V*eta)=1+C\_T*eta. \steptag{validated} \\[4pt]
\step{1.3.3} Since C\_theta,C\_V>=0 and C\_T=C\_theta+3*C\_V, C\_T>=C\_V. Therefore ||tilde-Delta||\_cb<=1+C\_V*eta<=1+C\_T*eta and the inverse-composition branch gives ||tilde-Upsilon||\_cb<=1+C\_T*eta. Taking their maximum proves the cb assertion, and the factor branch supplies the root's two-sided amplification assertion. \steptag{validated} \\[4pt]
\end{proofbox}

\end{document}
